\BourbakiExerciseGroup

\begin{enumerate}
\def\labelenumi{\arabic{enumi})}
\item
  Let E be an extension of a field K and let B be a transcendence basis of E over K. Show that E is equipotent to \(K \times B\) if one of the sets K, B is infinite, and countable otherwise. * Deduce in particular that every transcendence basis of the field R of real numbers over the field Q of rational numbers has the power of the continuum. *
\item
  Let K be the field \(\mathbf{Q}(\mathbf{X})\) of rational fractions in one indeterminate over the field Q of rational numbers. Show that in the ring K{[}Y{]} the polynomial \(Y^{2} + X^{2} + 1\) is irreducible and if E is the extension of K generated by a root of this polynomial (in an algebraic closure of
\end{enumerate}

K), then Q is relatively algebraically closed in E but E is not a pure extension of Q. (To see that there is in E no element \(a \notin Q\) algebraic over Q note that the existence of such an element would imply the relation \(\mathrm{E} = \mathbf{Q}(a)(\mathbf{X})\) , and observe that \(-(\mathbf{X}^{2} + 1)\) is not a square in \(\Omega(\mathbf{X})\) , where \(\Omega\) is an algebraic closure of Q.) Show that if i is a root of \(X^{2} + 1\) , then \(\mathrm{E}(i)\) is a pure transcendental extension of \(\mathbf{Q}(i)\) .

* 3) Let K be the field C(X) of rational fractions in one indeterminate over the (algebraically closed) field C of complex numbers. Show that in the ring K{[}Y{]}, the polynomial \(Y^{3} + X^{3} + 1\) is irreducible; let E be the extension of K generated by a root of this polynomial (in an algebraic closure of K). Show that E is not a pure extension of C, even though C is algebraically closed. (Show that it is impossible for a relation \(u^{3} + v^{3} + w^{3} = 0\) to exist between three polynomials u, v, w of C{[}X{]}, pairwise relatively prime and not all three constant. Argue by contradiction: if r is the largest of the degrees of u, v, w and if for example \(\deg(w) = r\) , deduce from the relation

\[
w ^ {3} = - (u + v) (u + j v) (u + j ^ {2} v)
\]

where \(j\) is a primitive cube root of unity, the existence of three polynomials \(u_1\) , \(v_1\) , \(w_1\) in \(\mathbf{C}[\mathbf{X}]\) , pairwise relatively prime, not all three constant, of degrees \(< r\) and such that \(u_1^3 + v_1^3 + w_1^3 = 0\) .) *

\begin{enumerate}
\def\labelenumi{\arabic{enumi})}
\setcounter{enumi}{3}
\tightlist
\item
  Let \(\Omega\) be an algebraically closed extension of a field \(\mathbf{K}\) , of infinite transcendence degree over \(\mathbf{K}\) .
\end{enumerate}

\begin{enumerate}
\def\labelenumi{\alph{enumi})}
\tightlist
\item
  Show that there exist an infinity of K-endomorphisms of \(\Omega\) having as images subfields of \(\Omega\) distinct from \(\Omega\) , and with respect to which \(\Omega\) may have a transcendence degree equal to any cardinal up to at most tr. \(\deg_{\mathbf{K}}\Omega\) . * In particular, there exist an infinity of distinct Q-isomorphisms of the field C of complex numbers onto subfields of C distinct from C. * b) Show that for every subextension E of \(\Omega\) such that tr. \(\deg_{\mathbf{K}}E < \text{tr} \cdot \deg_{\mathbf{K}}\Omega\) , every K-isomorphism of E onto a subfield of \(\Omega\) may be extended to a K-automorphism of \(\Omega\) . Give examples where tr. \(\deg_{\mathbf{K}}E = \text{tr} \cdot \deg_{\mathbf{K}}\Omega\) and where there exists a K-isomorphism of E onto a subfield of \(\Omega\) which cannot be extended to a K-isomorphism of any extension F of E, contained in \(\Omega\) and distinct from E, onto a subfield of \(\Omega\) .
\end{enumerate}

\begin{enumerate}
\def\labelenumi{\arabic{enumi})}
\setcounter{enumi}{4}
\item
  Let \(\Omega\) be an algebraically closed extension of a field \(\mathbf{K}\) . Show that for every subextension \(\mathbf{E}\) of \(\Omega\) , which is transcendental over \(\mathbf{K}\) , there exist an infinity of \(\mathbf{K}\) -isomorphisms of \(\mathbf{E}\) onto a subfield of \(\Omega\) (if \(x \in \mathbf{E}\) is transcendental over \(\mathbf{K}\) , consider a subextension \(\mathbf{F}\) of \(\mathbf{E}\) such that \(x\) is transcendental over \(\mathbf{F}\) and \(\mathbf{E}\) algebraic over \(\mathbf{F}(x)\) , and show that there exist an infinity of \(\mathbf{F}\) -isomorphisms of \(\mathbf{E}\) onto a subfield of \(\Omega\) ).
\item
  Let \(\Omega\) be an algebraically closed extension of a field \(K\) , and \(N\) a subextension of \(\Omega\) . Show that if \(N\) is quasi-Galois over \(K\) and if \(E\) and \(E'\) are two conjugate extensions of \(K\) contained in \(\Omega\) , then \(N(E)\) and \(N(E')\) are conjugate over \(K\) . Conversely, if \(N\) has this property and if \(\operatorname{tr} \cdot \deg_{K} N\) is finite and \(< \operatorname{tr} \cdot \deg_{K} \Omega\) , then it is necessarily algebraic and quasi-Galois over \(K\) .
\item
  Let E be an extension of a field K and \(\operatorname{Aut}_{\mathbf{K}}(\mathbf{E})\) the group of all K-automorphisms of E. a) For every finitely generated subextension L of E over K, \(\operatorname{Aut}_{\mathbf{L}}(\mathbf{E})\) is a subgroup of \(\operatorname{Aut}_{\mathbf{K}}(\mathbf{E})\) . Show that when L ranges over the set of subextensions of E finitely generated over K, the groups \(\operatorname{Aut}_{\mathbf{L}}(\mathbf{E})\) form a fundamental system of neighbourhoods of the neutral element of \(\mathrm{Aut}_{\mathbf{K}}(\mathbf{E})\) for a topology compatible with the group structure of \(\mathrm{Aut}_{\mathbf{K}}(\mathbf{E})\) and for which \(\mathrm{Aut}_{\mathbf{K}}(\mathbf{E})\) is a separated and totally disconnected group.
\end{enumerate}

\begin{enumerate}
\def\labelenumi{\alph{enumi})}
\setcounter{enumi}{1}
\tightlist
\item
  When E is equipped with the discrete topology, the topology of \(\operatorname{Aut}_{\mathbf{K}}(\mathbf{E})\) is the coarsest in which the mapping \((u,x)\to u(x)\) of \(\operatorname{Aut}_{\mathbf{K}}(\mathbf{E})\times\mathbf{E}\) into E is continuous.
\end{enumerate}

* c) The group \(\operatorname{Aut}_{\mathbf{Q}}(\mathbf{R})\) is reduced to the neutral element (V, p.~153, Ex. 2 of § 9). * \(d)\) Let \(\mathbf{K}_0 \supseteq \mathbf{K}\) be the subextension of E consisting of the elements invariant under all K-automorphisms of E. Show that for E to be algebraic over \(\mathbf{K}_0\) it is necessary and sufficient that \(\operatorname{Aut}_{\mathbf{K}}(\mathbf{E})\) should be compact.

\begin{enumerate}
\def\labelenumi{\alph{enumi})}
\setcounter{enumi}{4}
\tightlist
\item
  Assume that K is infinite. If \(E = K(X)\) , the field of fractions in one indeterminate over K, show that \(\operatorname{Aut}_{K}(E)\) is an infinite discrete group (V, p.~148, Ex. 11). If K is of characteristic 0, the subgroup \(\Gamma\) of \(\operatorname{Aut}_{K}(E)\) generated by the automorphism \(\sigma : X \mapsto X + 1\) is closed in \(\operatorname{Aut}_{K}(E)\) and distinct from the latter, but has the same field of invariants.
\end{enumerate}

\begin{enumerate}
\def\labelenumi{\arabic{enumi})}
\setcounter{enumi}{7}
\tightlist
\item
  Let \(E\) be an extension of a field \(K\) .
\end{enumerate}

\begin{enumerate}
\def\labelenumi{\alph{enumi})}
\item
  Let \(x\) be an element of \(E - K\) . Show that there exists a pure extension \(L \subset E\) of \(K\) such that \(E\) is algebraic over \(L\) and \(x \notin L\) (consider the set of subextensions \(L\) of \(E\) which are pure transcendental over \(K\) and such that \(x \notin L\) ).
\item
  Let \(x \in E\) be a transcendental element over \(K\) and \(p\) an integer \(\geqslant 2\) . Show that there exists a pure transcendental extension \(L \subset E\) of \(K\) such that \(E\) is algebraic over \(L\) , \(x \notin L\) and \(x^p \in L\) (same method).
\item
  Let F be a subextension of E and let \(x \in E\) be an algebraic element over F and \(P \in F[X]\) its minimal polynomial. Show that there exists a pure extension \(L \subset E\) of K such that E is algebraic over L and P is irreducible over the field \(F(L)\) .
\end{enumerate}

\begin{enumerate}
\def\labelenumi{\arabic{enumi})}
\setcounter{enumi}{8}
\tightlist
\item
  An extension E of a field K is called a Dedekind extension of K if for every subextension F of E, F is identical with the subfield of elements of E invariant under the group \(\mathrm{Aut}_{\mathrm{F}}(\mathrm{E})\) of all F-automorphisms of E.
\end{enumerate}

\begin{enumerate}
\def\labelenumi{\alph{enumi})}
\item
  For an algebraic extension \(\mathbf{E}\) of \(\mathbf{K}\) to be Dedekind over \(\mathbf{K}\) it is necessary and sufficient that \(\mathbf{E}\) should be a Galois extension of \(\mathbf{K}\) .
\item
  Suppose that K is of characteristic p \textgreater{} 0. Show that a Dedekind extension of K is necessarily algebraic over K (hence Galois). (Use Ex. 8, b).)
\item
  If E is a Dedekind extension of K and L a pure transcendental subextension of E, distinct from K and such that E is an algebraic extension of L, show that E has infinite degree over L.
\end{enumerate}

\begin{enumerate}
\def\labelenumi{\arabic{enumi})}
\setcounter{enumi}{9}
\item
  Let E be an extension of a field K such that the mapping \(F \mapsto \operatorname{Aut}_{F}(E)\) is a bijection of the set of subextensions of E onto the set of subgroups of \(\operatorname{Aut}_{K}(E)\) . Show that E is then a Galois extension of K of finite degree. (Using Ex. 9, c), show first that E is necessarily algebraic over K, then use Ex. 9, a).
\item
  Let \(\mathbf{E}\) be an extension of a field \(\mathbf{K}\) and \(\Omega\) an algebraically closed extension of \(\mathbf{E}\) . Show that the following conditions are equivalent:
\end{enumerate}

α) E is a p-radical extension of K;

\(\beta)\) for every extension \(F\) of \(K\) , \(E \otimes_{K} F\) has only a single prime ideal;

γ) for every extension F of K, any two composite extensions of E and F are isomorphic;

δ) \(E \otimes_{K} \Omega\) has only a single prime ideal ;

ε) any two composite extensions of E and Ω are isomorphic.

\begin{enumerate}
\def\labelenumi{\arabic{enumi})}
\setcounter{enumi}{11}
\tightlist
\item
  \begin{enumerate}
  \def\labelenumii{\alph{enumii})}
  \tightlist
  \item
    Let K be a field, \(\Omega\) an algebraically closed extension of K, \(E \subset \Omega\) and F two extensions of K; assume that \(\operatorname{tr} \cdot \deg_E \Omega \geqslant \operatorname{tr} \cdot \deg_K F\) . Show that every composite extension of E and F is isomorphic to a composite of the form \((L_w, 1, w)\) , where \(w\) is a K-isomorphism of F onto a subfield of \(\Omega\) and \(L_w\) is the subfield of \(\Omega\) generated by E and \(w(F)\) .
  \end{enumerate}
\end{enumerate}

\begin{enumerate}
\def\labelenumi{\alph{enumi})}
\setcounter{enumi}{1}
\tightlist
\item
  Deduce from a) that if F is algebraic of finite degree n over K, then there are at most n composite extensions of E and F that are pairwise non-isomorphic.
\end{enumerate}

\begin{enumerate}
\def\labelenumi{\arabic{enumi})}
\setcounter{enumi}{12}
\item
  In an algebraic closure \(\Omega\) of \(\mathbf{Q}(\mathbf{X})\) consider two pure transcendental extensions \(\mathrm{E} = \mathbf{Q}(\mathbf{X})\) and \(\mathrm{F} = \mathbf{Q}(\mathbf{X} + i)\) (where \(i^2 = -1\) ) of the field \(\mathbf{Q}\) . Show that \(\mathrm{E} \cap \mathrm{F} = \mathbf{Q}\) but that \(\mathrm{E}\) and \(\mathrm{F}\) are not algebraically disjoint over \(\mathbf{Q}\) .
\item
  Let E, F, G be three extensions of a field K contained in an extension \(\Omega\) of K and such that \(F \subset G\) . Show that for E and G to be algebraically disjoint over K it is necessary and sufficient that E and F should be algebraically disjoint over K and \(\mathrm{E}(\mathrm{F})\) and G algebraically disjoint over F.
\item
  \begin{enumerate}
  \def\labelenumii{\alph{enumii})}
  \tightlist
  \item
    Let K be a field and L a subfield of K such that K is a finitely generated L-algebra, in other words \(K = L[a_{1}, a_{2}, \ldots, a_{n}]\) for elements in K. Show that the \(a_{j}\) are algebraic over L. (Argue by contradiction: if \(a_{1}, \ldots, a_{m}\) ( \(m \geqslant 1\) ) form a maximal algebraically free subfamily of \((a_{j})_{1 \leqslant j \leqslant n}\) , observe that in the ring \(A = L[a_{1}, \ldots, a_{m}]\) the intersection of all non-zero ideals is reduced to zero, and use Ex. 5 of V, p.~147.)
  \end{enumerate}
\end{enumerate}

\begin{enumerate}
\def\labelenumi{\alph{enumi})}
\setcounter{enumi}{1}
\tightlist
\item
  Deduce from a) that in the polynomial algebra \(K[X_{1}, \ldots, X_{n}]\) every maximal ideal is of finite codimension over K. Conclude that if A is a commutative finitely generated algebra over K, then every subfield of A containing K is of finite degree over K.
\end{enumerate}

\begin{enumerate}
\def\labelenumi{\arabic{enumi})}
\setcounter{enumi}{15}
\tightlist
\item
  \begin{enumerate}
  \def\labelenumii{\alph{enumii})}
  \tightlist
  \item
    Let K be an algebraically closed field, \(\Omega\) an algebraically closed extension of K, E a subextension of \(\Omega\) and \(x\) an element of \(\Omega\) transcendental over E. Let \(\overline{\mathbf{K}(x)}\) be the relative algebraic closure of \(\mathbf{K}(x)\) in \(\Omega\) and let \(\mathbf{M} = \operatorname{E}(\overline{\mathbf{K}(x)})\) ; show that M is an algebraic extension of infinite degree of \(\operatorname{E}(x)\) . (Observe that for every integer \(m > 1\) the polynomial \(\mathbf{X}^m - x\) is irreducible in \(\operatorname{E}(x)[\mathbf{X}]\) (V, p.~91, Remark).
  \end{enumerate}
\end{enumerate}

\begin{enumerate}
\def\labelenumi{\alph{enumi})}
\setcounter{enumi}{1}
\tightlist
\item
  Let E be an algebraically closed field and F a finitely generated extension of E. Show that every algebraically closed subfield of F is necessarily contained in E. (Otherwise there would exist a subfield \(L \subset F\) which is algebraically closed and an element \(x \in L\) which is transcendental over E. Apply a), taking for K the relative algebraic closure in E of the prime subfield of E, noting that M should be finitely generated over E.)
\end{enumerate}

\begin{enumerate}
\def\labelenumi{\arabic{enumi})}
\setcounter{enumi}{16}
\item
  Let K be a field, \(\Omega\) an algebraically closed extension of K and x an element of \(\Omega\) transcendental over K. Show that for every prime number q there exists a subextension M of \(\Omega\) such that \(\mathbf{M}(x)\) is a Galois extension of degree q of M. (Consider the element \(y = x^{q}\) if q is not the characteristic of K, \(y = x^{q} - x\) in the opposite case, and a subextension L of \(\Omega\) generated by a transcendence basis of \(\Omega\) over K containing y and besides the q-th roots of unity if \(y = x^{q}\) ; use Ex. 12 of V, p.~163.)
\item
  Let K be a field and \(\mathbf{K}((\mathbf{X}))\) the field of formal power series in one indeterminate over K (IV, p.~38). Show that \(\operatorname{tr}.\deg_{\mathbf{K}}\mathbf{K}((\mathbf{X}))=\operatorname{Card}(\mathbf{K}^{\mathbf{N}})\) . Distinguish two cases:
\end{enumerate}

\begin{enumerate}
\def\labelenumi{\alph{enumi})}
\item
  If \(\operatorname{Card}(\mathbf{K}) < \operatorname{Card}(\mathbf{K}^{\mathbf{N}})\) , note that \(\operatorname{Card}(\mathbf{K}((\mathbf{X}))) = \operatorname{Card}(\mathbf{K}^{\mathbf{N}})\) .
\item
  If \(\operatorname{Card}(K) = \operatorname{Card}(K^{N})\) , let P be the prime subfield of K, S an infinite set of elements of \(\operatorname{K}((X))\) which is algebraically independent over K and T the set of coefficients of all the formal power series belonging to S. Let L be the relative algebraic closure in
\end{enumerate}

K((X)) of the field K(S) and u an element of L; there is an equation \(g(s_{1}, \ldots, s_{m}, u) = 0\) where \(g \neq 0\) is a polynomial in \(K[X_{1}, \ldots, X_{m}, X_{m+1}]\) and \(s_{1}, \ldots, s_{m}\) elements of S. Let A be the set of coefficients of the polynomial g and \(C(u)\) the set of coefficients of the formal power series u; show that the field \(P(T)(A \cup C(u))\) is algebraic over \(P(T \cup A)\) (denoting by \(\Omega\) an algebraic closure of K, show that otherwise there would exist an infinity of formal power series \(v \in \Omega((X))\) satisfying the equation g ( \(s_{1}, \ldots, s_{m}, v) = 0\) ). Using Ex. 1, show that if \(\text{Card}(S) < \text{Card}(K)\) , then the transcendence degree of K over \(P(T)\) is infinite and deduce that in this case L is distinct from \(K((X))\) (if \((t_{n})_{n \geqslant 0}\) is an infinite sequence of elements of K algebraically independent over \(P(T)\) , consider the formal power series \(\sum_{n=0}^{\infty} t_{n} X^{n}\) ).

\begin{enumerate}
\def\labelenumi{\arabic{enumi})}
\setcounter{enumi}{18}
\tightlist
\item
  Let \(\mathbf{E}\) and \(\mathbf{F}\) be two transcendental extensions of a field \(\mathbf{K}\) , linearly disjoint over \(\mathbf{K}\) . Show that \(\mathbf{K}(\mathbf{E} \cup \mathbf{F})\) is distinct from the ring \(\mathbf{C}\) (isomorphic to \(\mathbf{E} \otimes_{\mathbf{K}} \mathbf{F}\) ) generated by \(\mathbf{E} \cup \mathbf{F}\) . (Reduce to the case where the transcendence degrees of \(\mathbf{E}\) and \(\mathbf{F}\) over \(\mathbf{K}\) are equal to 1; if \(x \in \mathbf{E}\) and \(y \in \mathbf{F}\) are transcendental over \(\mathbf{K}\) show that we cannot have \((x + y)^{-1} \in \mathbf{C}\) ; arguing by contradiction, show that in the opposite case, if \(p\) is the characteristic exponent of \(\mathbf{K}\) , there would exist an integer \(r \geqslant 0\) such that \((x + y)^{-p^r}\) belonged to the subring of \(\mathbf{C}\) generated by \(\mathbf{K}(x) \cup \mathbf{K}(y)\) .)
\end{enumerate}

* 20) Let \(n\) be an integer, \(K\) a field of characteristic \(p\) prime to \(n\) , possessing primitive \(n\) -th roots of unity, \(\mu_n \subset K^*\) the subgroup of \(n\) -th roots of unity, \(G\) a finite commutative group annihilated by \(n\) and \(A = K^{(G)}\) the group algebra of \(G\) (III, p.~446).

\begin{enumerate}
\def\labelenumi{\alph{enumi})}
\item
  Show that A is a diagonalizable algebra (V, p.~28). (Decompose G into a direct sum of cyclic groups (VII, p.~22).) Deduce that every A-module is a direct sum of modules which are vector K-spaces of dimension 1.
\item
  Let V be an A-module of finite dimension over K, S the symmetric algebra of the vector K-space V and F the field of fractions of S. The group G operates on V, hence on S and F. Assume that G acts faithfully on V. Show that G operates faithfully on F. Calculate \([F : F^{G}]\) .
\item
  Let \((x_{1},\ldots ,x_{n})\) be a basis of the vector K-space V consisting of eigenvectors for the action of G. Let \(\Gamma\) be the subgroup of \(\mathbf{F}^*\) generated by \(x_{1},\ldots ,x_{n}\) . Show that \(\Gamma\) is a free commutative group and that the sub-K-algebra B of F generated by \(\Gamma\) may be identified with \(\mathbf{K}^{(\Gamma)}\) and is stable under the operation of G.
\item
  For all \(x \in \Gamma\) , the mapping \(\gamma \mapsto \frac{\gamma x}{x}\) is a homomorphism of \(G\) into \(\mu_n\) , giving rise to a homomorphism \(\theta: \Gamma \to \operatorname{Hom}(G, \mu_n)\) . Show that \(\theta\) is surjective. Put \(\Gamma_1 = \operatorname{Ker}(\theta)\) . Using the elementary divisor theorem (VII, p.~24), show that \(B\) is a free \(K^{(\Gamma_1)}\) -module of rank Card(G).
\item
  By passing to the field of fractions deduce from \(d\) and \(a\) that \(\mathbf{F}^{\mathrm{G}}\) is the field of fractions of \(\mathbf{K}^{(\Gamma_1)}\) . Deduce that \(\mathbf{F}^{\mathrm{G}}\) is a pure transcendental extension of \(\mathbf{K}\) . *
\end{enumerate}

